\section{Discussion and conclusion}
\label{sec:discussion}

\textbf{The horizon of the reward selects the reward hack.} Which exploit a policy finds depends
on what the judge is shown. The $K{=}0$ arm finds unearned praise, whose payoff is immediate and
whose cost is deferred. The $K{=}5$ arm cannot profit from that, because
the cost now lands inside the window the judge sees, and it moves toward asking questions and, to
a smaller degree, toward pushing. Extending the horizon did not make the policy honest; it made
the dishonest strategy unprofitable.

\textbf{Scope: one optimiser, one regime.} Everything here concerns a group-relative optimiser.
The lever's origin is a preference-pair method: \citet{baruch2025pto} introduced it with
preference trees and DPO in an earlier, easier regime of this task (a 7B policy, more cooperative
patients, a weaker judge). Whether the lever helps other optimiser
families in other regimes is outside this paper's scope. What would change our minds: a second
\emph{training} run reversing the ordering, an intermediate $K$ showing a threshold rather than a
horizon, or a human MI coder disagreeing with the over-praise finding.

\textbf{Ceiling-driven grader saturation is worth naming.} Optimising against an LLM judge
succeeds first on the judge's own terms: the scores it gives move to the top of its scale, and a
bounded scale at its top cannot rank. Here the trained-against grader lost three-quarters of its
per-conversation variance while a harsher independent grader's spread stayed flat, at the
checkpoint a practitioner would ship, and arm-level agreement was blind to it. The diagnostic is
cheap: track the training grader's per-conversation variance and share of ceiling scores, not
only its mean, and keep a second grader with headroom.

\textbf{Conclusion.} For a group-relative optimiser in multi-turn dialogue, what the judge
is shown is a first-order design choice. Scoring the continuation rather than the turn raised what
the same optimiser extracted from the same oracle, under a held-out grader, against the turn-level
arm's best checkpoint and on a replicated draw, and decided which failure mode the policy
learned, with no new loss, value function, or tree. The reward horizon
deserves to be reported as a hyperparameter, and the training grader's variance to be watched as
a statistic.
